\scanpage{52}
\handwritten{%
réalisation dans $E$ sont variables\ldots)

b) Polygône régulier. On avait proposé \inserted{\unsure{la|de}} (dans contexte
métrique, \inserted{sur $\mathbb R$}) : les \inserted{groupe $Z$ des} \unsure{déplacement|déplacements} qui invarient le
polygône doivent être transitifs sur l'ens des
sommets. \deleted{Dans} \inserted{Nous restreignant pour simplifier} au cas fini \unsure{Ceci|Ceci il|il} est clair :
\unsure{primo|priori} qu'il y a un pt invariant -- centre de
gravité des l'ens des sommets -- \deleted{le groupe des} \inserted{les \uline{déplacements}}
\inserted{sont des} rotations, \deleted{et le groupe} le groupe \inserted{$R$} des rotations
de centre $o$ étant le ``groupe des angles'', can\textsuperscript{t}
isom. à $\mathbb R/\mathbb Z\simeq U=\{z\in\mathbb C\mid z\bar z\unsure{)|\rvert}=1\}$. On
trouvait, si $\zeta\in\unsure{R|\mathbb R}$ \deleted{est} tel que
\[
\zeta s_0=s_1
\]
\[
\text{que }\left\{\begin{aligned}
&\zeta^n=1\qquad\zeta^d\ne1\quad\text{si }1\leq d<n,\ \text{et les}\\
&S=\{\underbrace{\zeta^i s_0}_{s_i}\mid i\in\mathbb Z/n\mathbb Z\}\\
&a_i=\mathrm{Dr}(s_i,s_{i+1})
\end{aligned}\right.
\]
% slanted
\margin{en fait, $S$ et $A$ sont des \uline{torseurs} sous $Z$ i.e. $Z$ est simplement et transitif sur $S$ et $A$}
et inversement, il est clair que les polygônes
ainsi obtenus sont ``réguliers'' en le sens proposé.

Si on prend les autom \uline{euclidiens} de $E$
sans plus \inserted{qui \unsure{conservent|stabilisent|laissent} $P$}, ils ont une propriété \deleted{plus} de transitivité
plus forte : \deleted{ils sont} \inserted{leur groupe $D$ est} \unsure{transitifs|transitif} sur l'ens
des \uline{repères} du polygône (et évidemment, = simplement
transitif\unsure{\}|)}, puisque $S$ engendre affinement $E$,
}
% letterhead
% Printed upside down relative to the notes.
\begin{tabular}{@{}l@{\qquad}l@{}}
UNIVERSITE DES SCIENCES ET TECHNIQUES &\\
DU LANGUEDOC &\\
\rule{1.4em}{.35pt}& Montpellier, le\\
INSTITUT DE MATHEMATIQUES &\\
Place Eugène Bataillon &\\
34060 MONTPELLIER CEDEX &\\
\rule{1.4em}{.35pt}&
\end{tabular}
% archive sheet number: 51
